\documentclass[12pt,a4paper]{article}

\usepackage[utf8]{inputenc}
\usepackage[T1]{fontenc}
\usepackage[spanish,es-nodecimaldot]{babel}
\usepackage[top=2.5cm,bottom=2.5cm,left=3cm,right=2.5cm]{geometry}
\usepackage{xcolor}
\usepackage{tcolorbox}
\usepackage{enumitem}
\usepackage{array}
\usepackage{booktabs}
\usepackage{fancyhdr}
\usepackage{titlesec}
\usepackage{parskip}
\usepackage{multirow}
\usepackage{tabularx}
\usepackage{hyperref}

% ── Colores ─────────────────────────────────────────────────────────────────
\definecolor{azulSMR}  {RGB}{0,90,170}
\definecolor{grisCard} {RGB}{245,247,250}
\definecolor{grisOsc}  {RGB}{80,80,100}
\definecolor{rojoEntr} {RGB}{180,30,30}
\definecolor{verdeOK}  {RGB}{0,120,70}
\definecolor{ambarInfo}{RGB}{145,95,0}

% ── Cajas tcolorbox ──────────────────────────────────────────────────────────
\tcbuselibrary{skins,breakable}

\newtcolorbox{contextobox}{
  enhanced, breakable,
  colback=azulSMR!6, colframe=azulSMR,
  arc=3pt, boxrule=0.8pt,
  fonttitle=\bfseries\color{white},
  coltitle=white, attach boxed title to top left={yshift=-2mm,xshift=4mm},
  boxed title style={colback=azulSMR, arc=2pt, boxrule=0pt},
  title=Contexto del proyecto
}

\newtcolorbox{tablabox}[1]{
  enhanced, breakable,
  colback=grisCard, colframe=grisOsc!60,
  arc=3pt, boxrule=0.7pt,
  fonttitle=\bfseries,
  title=#1,
  left=4pt, right=4pt, top=4pt, bottom=4pt
}

\newtcolorbox{consultabox}{
  enhanced, breakable,
  colback=verdeOK!5, colframe=verdeOK,
  arc=3pt, boxrule=0.8pt,
  fonttitle=\bfseries\color{white},
  coltitle=white, attach boxed title to top left={yshift=-2mm,xshift=4mm},
  boxed title style={colback=verdeOK, arc=2pt, boxrule=0pt},
  title=Bloque 3 — Consultas SQL requeridas
}

\newtcolorbox{entregabox}{
  enhanced,
  colback=rojoEntr!6, colframe=rojoEntr,
  arc=3pt, boxrule=1pt,
  fonttitle=\bfseries\color{white},
  coltitle=white, attach boxed title to top left={yshift=-2mm,xshift=4mm},
  boxed title style={colback=rojoEntr, arc=2pt, boxrule=0pt},
  title=ENTREGABLE — Qué debes presentar
}

\newtcolorbox{notabox}{
  colback=ambarInfo!8, colframe=ambarInfo,
  arc=3pt, boxrule=0.7pt,
  fonttitle=\bfseries\color{ambarInfo},
  title=Nota del profesor
}

% ── Cabecera/pie ─────────────────────────────────────────────────────────────
\pagestyle{fancy}
\fancyhf{}
\fancyhead[L]{\small\color{grisOsc} Módulo AOF — Bases de Datos con Access}
\fancyhead[R]{\small\color{grisOsc} 1.º SMR · IES\,---}
\fancyfoot[C]{\small\thepage}
\fancyfoot[R]{\small\color{grisOsc} Práctica: TechFix SL}
\renewcommand{\headrulewidth}{0.4pt}
\renewcommand{\footrulewidth}{0.2pt}

% ── Secciones ────────────────────────────────────────────────────────────────
\titleformat{\section}
  {\large\bfseries\color{azulSMR}}{}{0em}{\thesection.\ }[\vspace{-4pt}\color{azulSMR}\rule{\linewidth}{0.6pt}]
\titleformat{\subsection}
  {\normalsize\bfseries\color{grisOsc}}{}{0em}{\thesubsection\ }

\setlist[itemize]{leftmargin=1.6em, itemsep=1pt, topsep=3pt}
\setlist[enumerate]{leftmargin=1.6em, itemsep=3pt, topsep=3pt}

\hypersetup{colorlinks=true, linkcolor=azulSMR, urlcolor=azulSMR, pdfborder={0 0 0}}

% ═══════════════════════════════════════════════════════════════════════════════
\begin{document}
% ═══════════════════════════════════════════════════════════════════════════════

% ── PORTADA ──────────────────────────────────────────────────────────────────
\begin{center}
  {\color{azulSMR}\rule{\linewidth}{2.5pt}}\\[10pt]
  {\Huge\bfseries\color{azulSMR} PRÁCTICA INTEGRADORA}\\[4pt]
  {\LARGE\bfseries Bases de Datos con Microsoft Access}\\[6pt]
  {\large\color{grisOsc} Diseño · Implementación · Consultas SQL}\\[10pt]
  {\color{azulSMR}\rule{\linewidth}{1pt}}\\[8pt]
  \renewcommand{\arraystretch}{1.3}
  \begin{tabular}{@{}ll}
    \bfseries Módulo:       & Aplicaciones Ofimáticas (AOF) \\
    \bfseries Ciclo:        & Técnico en Sistemas Microinformáticos y Redes (SMR) \\
    \bfseries Curso:        & 1.º \\
    \bfseries Unidad:       & UT6 — Utilización de Bases de Datos Ofimáticas \\
    \bfseries Modalidad:    & Individual, entregable \\
    \bfseries Calificación: & Hasta 10 puntos (ver rúbrica) \\
  \end{tabular}\\[8pt]
  \colorbox{azulSMR!15}{\parbox{0.85\textwidth}{\centering\small
    Rellena los datos antes de entregar:\\[4pt]
    \textbf{Nombre y apellidos:}\enspace\underline{\hspace{7cm}}
    \quad \textbf{Grupo:}\enspace\underline{\hspace{1.5cm}}
  }}
\end{center}

\vspace{0.4cm}

% ── 1. CONTEXTO ──────────────────────────────────────────────────────────────
\section{Contexto del proyecto}

\begin{contextobox}
\textbf{TechFix SL} es un pequeño taller de reparación de equipos informáticos situado en el centro
de la ciudad. El negocio atiende a particulares y pequeñas empresas y necesita digitalizar su gestión.
Actualmente lleva el control en papel y quiere pasar a una base de datos que le permita:

\begin{itemize}
  \item Registrar los clientes y los dispositivos que traen a reparar.
  \item Asignar órdenes de trabajo a los técnicos del taller.
  \item Controlar el inventario de repuestos y vincularlos a cada reparación.
  \item Consultar el historial de reparaciones y generar informes.
\end{itemize}

\textbf{Tu misión} es diseñar e implementar esta base de datos completa en Microsoft Access,
redactar la documentación técnica y escribir las consultas SQL que el taller necesita.
\end{contextobox}

\vspace{0.3cm}

% ── 2. DISEÑO ────────────────────────────────────────────────────────────────
\section{Bloque 1 — Diseño de la base de datos}

\subsection{Requisitos de las entidades}

Debes identificar las tablas necesarias y definir sus campos, tipos de dato y claves.
A continuación se describen los \textbf{requisitos funcionales} de cada entidad.
\textbf{No se especifican los nombres exactos de los campos}: tú decides el nombre, el tipo
y si es obligatorio u opcional.

\vspace{0.3cm}

\begin{tablabox}{Entidad 1 — Clientes}
\begin{itemize}
  \item Cada cliente debe tener un identificador único generado automáticamente.
  \item Debe almacenarse el nombre completo (nombre y apellidos).
  \item Se necesitan \textbf{al menos dos formas de contacto} (p.\ ej.\ teléfono y correo).
  \item Se debe registrar la localización: como mínimo localidad y provincia.
  \item Ha de quedar constancia de cuándo se registró por primera vez en el taller.
  \item Debe ser posible saber si el cliente está activo (acepta comunicaciones) o no.
\end{itemize}
\end{tablabox}

\vspace{0.25cm}

\begin{tablabox}{Entidad 2 — Técnicos}
\begin{itemize}
  \item Identificador único y datos personales básicos.
  \item Cada técnico tiene \textbf{una única especialización principal}: hardware, software o redes.
  \item Debe reflejarse si está activo o de baja en la empresa.
  \item Se registra la fecha en que se incorporó al taller.
  \item No es necesario almacenar contraseñas ni accesos al sistema.
\end{itemize}
\end{tablabox}

\vspace{0.25cm}

\begin{tablabox}{Entidad 3 — Dispositivos (equipos)}
\begin{itemize}
  \item Cada dispositivo pertenece a \textbf{un único cliente} (un cliente puede tener varios).
  \item Se clasifica por tipo: ordenador de sobremesa, portátil, tablet, smartphone, impresora u otro.
  \item Se almacenan la marca y el modelo del fabricante.
  \item El cliente describe brevemente el problema al traer el equipo (campo de texto libre).
  \item Un mismo dispositivo puede traerse a reparar en distintas ocasiones.
\end{itemize}
\end{tablabox}

\vspace{0.25cm}

\begin{tablabox}{Entidad 4 — Órdenes de trabajo (tabla central)}
\begin{itemize}
  \item Es la entidad central del sistema: vincula un dispositivo con un técnico asignado.
  \item Debe registrarse la fecha y hora de entrada del equipo al taller.
  \item Se indica una fecha estimada de entrega al cliente.
  \item El estado de la orden debe ser uno de los siguientes: \textit{Recibido, En diagnóstico,
        En reparación, En espera de piezas, Terminado, Entregado}.
  \item El técnico debe poder añadir notas de diagnóstico y descripción del trabajo realizado
        (pueden ser dos campos distintos).
  \item Se registra el coste de la mano de obra de forma independiente al coste de los repuestos.
  \item Una orden de trabajo puede no tener técnico asignado inicialmente.
\end{itemize}
\end{tablabox}

\vspace{0.25cm}

\begin{tablabox}{Entidad 5 — Repuestos}
\begin{itemize}
  \item Catálogo de piezas disponibles en el taller.
  \item Cada repuesto tiene un nombre, una referencia interna y una categoría
        (RAM, disco duro/SSD, pantalla, teclado, batería, fuente de alimentación, otro).
  \item Se distingue entre el precio de coste (lo que paga el taller) y el precio de venta al cliente.
  \item Se controla el stock: número de unidades disponibles en almacén.
  \item Incluye el nombre del proveedor habitual.
\end{itemize}
\end{tablabox}

\vspace{0.25cm}

\begin{tablabox}{Entidad 6 — Uso de repuestos en órdenes (tabla intermedia)}
\begin{itemize}
  \item Registra qué repuestos se han utilizado en cada orden de trabajo.
  \item Una misma orden puede usar varios repuestos distintos; un mismo repuesto puede
        aparecer en múltiples órdenes: relación \textbf{N:M}.
  \item Se almacena la cantidad de unidades del repuesto usadas en esa orden.
  \item Se guarda el precio de venta aplicado en el momento de la reparación (puede diferir
        del precio actual del catálogo).
\end{itemize}
\end{tablabox}

\vspace{0.4cm}

\subsection{Diagrama Entidad-Relación}

Antes de crear las tablas en Access, elabora el \textbf{diagrama E-R} de la base de datos.
Puedes dibujarlo a mano en papel cuadriculado o con cualquier herramienta (draw.io, Lucidchart,
papel y escáner\ldots). Debe mostrar:

\begin{itemize}
  \item Todas las entidades y sus atributos principales.
  \item Todas las relaciones entre entidades con su cardinalidad (1:1, 1:N, N:M).
  \item Indicación de claves primarias y claves foráneas.
\end{itemize}

\vspace{0.4cm}

% ── 3. IMPLEMENTACIÓN ────────────────────────────────────────────────────────
\section{Bloque 2 — Implementación en Microsoft Access}

\begin{enumerate}[label=\bfseries\arabic{enumi}.]
  \item \textbf{Crear la base de datos} en Access con el nombre \texttt{TechFix\_Apellido.accdb}.
  \item \textbf{Crear todas las tablas} con sus campos, tipos de datos y propiedades
        (tamaño, formato, valor predeterminado, requerido\ldots).
  \item \textbf{Definir las claves primarias} en cada tabla.
  \item \textbf{Establecer las relaciones} entre tablas en la ventana de relaciones de Access,
        activando la \textbf{integridad referencial} en todas ellas.
  \item \textbf{Introducir datos de prueba}: al menos 5 clientes, 3 técnicos, 8 dispositivos,
        10 órdenes de trabajo en distintos estados, y 6 repuestos con usos en varias órdenes.
  \item \textbf{Crear como mínimo los siguientes elementos de interfaz}:
        \begin{itemize}
          \item \textbf{Formulario de alta de cliente} (con todos sus campos).
          \item \textbf{Formulario de nueva orden de trabajo} (con subformulario de repuestos usados).
          \item \textbf{Informe de órdenes abiertas} agrupado por técnico, con nombre del cliente
                y dispositivo.
          \item \textbf{Informe de factura} para una orden concreta: muestra datos del cliente,
                dispositivo, trabajo realizado, repuestos y coste total.
        \end{itemize}
\end{enumerate}

\vspace{0.4cm}

% ── 4. CONSULTAS ─────────────────────────────────────────────────────────────
\section{Bloque 3 — Consultas SQL}

\begin{consultabox}
Crea en Access las siguientes consultas en \textbf{vista SQL}. Guarda cada una con el nombre
indicado y haz una captura de pantalla con los resultados usando tus datos de prueba.

\vspace{0.3cm}
\begin{enumerate}[label=\bfseries C\arabic*.]

  \item \textbf{[C1\_DispositivosPorLocalidad]} —
        Lista el nombre del cliente, tipo de dispositivo y marca de todos los equipos
        traídos por clientes de una localidad concreta (pon la tuya o invéntate una).
        Muestra solo aquellos cuya orden tenga estado distinto de \textit{Entregado}.

  \item \textbf{[C2\_OrdenesAbiertas]} —
        Muestra todas las órdenes de trabajo que no están en estado \textit{Entregado} ni
        \textit{Terminado}, ordenadas de más antigua a más reciente por fecha de entrada.
        Incluye el nombre del técnico asignado (puede ser \texttt{NULL} si no hay ninguno).

  \item \textbf{[C3\_RepPorTecnico]} —
        Cuenta el número de reparaciones completadas (estado \textit{Terminado} o
        \textit{Entregado}) que ha realizado cada técnico.
        Muestra solo los técnicos con \textbf{más de 1} reparación completada,
        ordenados de mayor a menor número.

  \item \textbf{[C4\_IngresosPorOrden]} —
        Calcula el coste total de cada orden de trabajo sumando la mano de obra más
        el importe de todos los repuestos usados (cantidad $\times$ precio aplicado).
        Muestra solo órdenes en estado \textit{Terminado} o \textit{Entregado},
        ordenadas por coste total descendente.

  \item \textbf{[C5\_RepuestoMasUsado]} —
        Obtén el nombre del repuesto que más veces ha aparecido en órdenes de trabajo
        (mayor número de órdenes distintas en las que se ha utilizado).

  \item \textbf{[C6\_ClientesVariosDispositivos]} —
        Lista los clientes que han traído \textbf{más de un dispositivo} distinto al taller,
        mostrando nombre completo y número de dispositivos.

  \item \textbf{[C7\_StockBajo]} —
        Muestra todos los repuestos cuyo stock sea igual o inferior a 3 unidades,
        ordenados por cantidad disponible de menor a mayor.

\end{enumerate}
\end{consultabox}

\vspace{0.4cm}

% ── 5. ENTREGABLE ────────────────────────────────────────────────────────────
\section{Entregable y criterios de evaluación}

\begin{entregabox}
Comprime todo en un único archivo \texttt{ZIP} con el nombre \texttt{TechFix\_Apellido.zip}
y entrégalo en la plataforma antes de la fecha indicada por el profesor. Debe contener:

\begin{enumerate}
  \item \textbf{Archivo de base de datos} \texttt{TechFix\_Apellido.accdb} con datos de prueba.
  \item \textbf{Documento de diseño} (PDF o Word) con:
    \begin{itemize}
      \item Diagrama E-R (foto o imagen del dibujo a mano también es válido).
      \item Descripción de cada tabla: nombre, campos elegidos, tipos de datos y justificación
            de las decisiones tomadas.
      \item Diagrama de relaciones (captura de la ventana de relaciones de Access).
    \end{itemize}
  \item \textbf{Capturas de pantalla} de los resultados de cada una de las 7 consultas SQL
        con tus datos de prueba (una imagen por consulta, con el nombre de la consulta visible).
  \item \textbf{Capturas de pantalla} de los dos formularios y los dos informes en funcionamiento.
\end{enumerate}
\end{entregabox}

\vspace{0.4cm}

% ── 6. RÚBRICA ───────────────────────────────────────────────────────────────
\section{Rúbrica de evaluación}

\renewcommand{\arraystretch}{1.4}
\begin{tabularx}{\linewidth}{@{}X c c c c@{}}
  \toprule
  \textbf{Criterio} & \textbf{Insuf.} & \textbf{Suf.} & \textbf{Bien} & \textbf{Sobres.} \\
   & \small(0--4) & \small(5--6) & \small(7--8) & \small(9--10) \\
  \midrule
  Diagrama E-R correcto y completo
    & Falta o muy incorrecto & Entidades y relaciones básicas & Cardinalidades correctas
    & Todo correcto + justificación \\
  \addlinespace
  Tablas, campos y tipos correctos (Access)
    & $<$4 tablas o tipos erróneos & Tablas básicas correctas & Todos los campos y tipos
    & Propiedades y validaciones extras \\
  \addlinespace
  Relaciones e integridad referencial
    & Sin relaciones & Algunas relaciones & Todas correctas & IR activada en todas \\
  \addlinespace
  Datos de prueba coherentes
    & Sin datos o incoherentes & Datos mínimos & Datos suficientes & Datos ricos y realistas \\
  \addlinespace
  Consultas SQL (C1--C7)
    & 0--2 funcionan & 3--4 funcionan & 5--6 funcionan & Las 7 funcionan correctamente \\
  \addlinespace
  Formularios e informes
    & No entregados & 1 formulario & 2 formularios + 1 informe & Todo completo y con formato \\
  \addlinespace
  Documento de diseño
    & No entregado & Solo capturas & Descripción básica & Justificación de decisiones \\
  \bottomrule
\end{tabularx}

\vspace{0.5cm}

\begin{notabox}
No existe una única solución correcta para el diseño de la base de datos. Se valorará la
\textbf{coherencia} de las decisiones tomadas y que el alumno sea capaz de \textbf{justificarlas}.
Por ejemplo: puede haber más o menos tablas que las descritas si existe una razón técnica válida.
Lo importante es que el modelo sea correcto, normalizado y funcional.

\medskip
Copia, plagio o entrega de trabajos ajenos supone \textbf{calificación de 0} en la práctica
y anotación en el expediente del alumno conforme al reglamento del centro.
\end{notabox}

\end{document}
